% =============================================================================
% MIVA-KNIGHT — Revised IT Project Proposal (aligned with repository docs)
% Supersedes narrative inconsistencies found across capstone .tex sources.
% Student context: IT 581 Capstone — Concordia University of Edmonton
% =============================================================================
\documentclass[12pt,a4paper]{article}

\usepackage[utf8]{inputenc}
\usepackage[T1]{fontenc}
\usepackage[margin=2.5cm]{geometry}
\usepackage{amsmath,amssymb}
\usepackage{booktabs}
\usepackage{enumitem}
\usepackage{hyperref}
\usepackage{xcolor}
\usepackage{setspace}
\setstretch{1.12}

\hypersetup{colorlinks=true,linkcolor=blue!50!black,urlcolor=blue!60!black}

\title{\textbf{MIVA-KNIGHT}\\[0.35em]
\large Revised Project Proposal:\\
Multimodal Perception, Shared Embeddings, and Grounded Voice QA}
\author{IT 581 Capstone (revised synthesis)}
\date{\today}

\begin{document}
\maketitle

\begin{abstract}
This document replaces earlier proposal wording that mixed \emph{runtime system stages}
with \emph{offline training tracks}, understated the number of experimental pipelines,
or implied a single monolithic multimodal model where the implementation is in fact a
\emph{composition} of frozen encoders, trainable projection heads, retrieval (FAISS),
knowledge-graph re-ranking, and a separate generator (LLM). The goal remains a
research prototype of a domain-adaptive, multimodal-aware voice assistant with
verifiable retrieval grounding; the revision states scope, dependencies, and
limitations in terms that match the technical documentation already produced for
pipelines A--E and the hybrid RAG stack.
\end{abstract}

\section{Problem and objective}

Users expect voice interfaces to use \textbf{text}, \textbf{speech}, and sometimes
\textbf{images} coherently. Large vision--language models offer one path, but this
project explores a modular architecture: specialised encoders map modalities into a
shared metric space, retrieval supplies evidence, and a language model synthesises
answers constrained by retrieved context. A secondary objective is \textbf{affect-aware
routing}: audio (and where available, video) cues inform emotion or sentiment estimates
that can modulate prompts or downstream behaviour.

\paragraph{Success criteria (engineering).}
\begin{itemize}[nosep]
  \item Demonstrate \textbf{cross-modal alignment} (e.g., image--text InfoNCE) into a
        common $512$-d unit hypersphere for retrieval.
  \item Demonstrate \textbf{domain transfer} for the visual branch (e.g., general
        images $\to$ radiology captions) with minimal retraining of the image-to-space
        adapter, as documented for the ROCO track.
  \item Demonstrate \textbf{grounded QA}: answers tied to retrieved documents with
        explicit confidence handling, not unconstrained hallucination.
  \item Demonstrate \textbf{standalone affect pipelines} where joint training with
        unrelated image--text pairs would corrupt the embedding space (lessons from
        Pipeline~C).
\end{itemize}

\section{System scope: two different ``phase'' vocabularies}

Earlier materials used ``Phase~1--4'' for two incompatible ideas. This proposal
\emph{separates} them:

\subsection*{Runtime stages (online assistant loop)}
The \textbf{hybrid RAG} stack documented in the capstone exposition is best described
as four \textbf{runtime stages}, not ``pipelines'':
\begin{enumerate}[label=(\arabic*)]
  \item \textbf{Encode} user query (text and/or audio via ASR/embeddings) into the
        shared retrieval embedding space.
  \item \textbf{Retrieve} top-$K$ candidates with FAISS, then re-rank with PMI-style
        knowledge-graph signals.
  \item \textbf{Generate} an answer from an LLM conditioned on retrieved context,
        with tiered confidence / abstention where appropriate.
  \item \textbf{Speak} (optional TTS) the response.
\end{enumerate}

\subsection*{Offline training tracks (research pipelines A--E)}
\textbf{Pipelines} are \emph{offline} experiments that train or freeze components.
The repository documentation describes \textbf{five} such tracks, not four:
\begin{center}
\small
\begin{tabular}{@{}llp{6.8cm}@{}}
\toprule
\textbf{Track} & \textbf{Primary data} & \textbf{Role in the larger system} \\
\midrule
A / general IR & COCO (text + image) &
  Image--text alignment; general-domain FAISS index and KG. \\
B / medical IR & ROCO (radiology text + image) &
  Medical-domain index; surgical transfer of the image projection head. \\
C / audio emotion & RAVDESS &
  Audio-only SupCon; emotion centroids; fixes prior joint-training contamination. \\
D / audio--video emotion & CREMA-D &
  Cross-modal InfoNCE + SupCon phases; inherits audio head warm-start from C. \\
E / multimodal sentiment & CMU-MOSI &
  Optional research branch: three-stream fusion (text + audio + video frame) for
  sentiment regression; uses pseudo-labels where gold labels are unavailable---reported
  metrics must be interpreted accordingly. \\
\bottomrule
\end{tabular}
\end{center}

\noindent
\textbf{Correction to prior text:} calling MIVA-KNIGHT a ``four-pipeline'' system
while also documenting Pipeline~E is internally inconsistent. The accurate statement
is: \emph{four core retrieval/affect tracks (A--D) plus one optional MOSI sentiment
fusion track (E).}

\section{Multimodal alignment assessment}

From a software and ML engineering perspective, the work \textbf{does align} with the
intent of building a multimodal system, with these clarifications:

\begin{itemize}
  \item \textbf{Modality coverage:} text, image, audio, and (in Pipeline~D/E) video
        frames are all represented. There is no single transformer that ingests raw
        pixels, waveforms, and tokens jointly in all paths; instead, \textbf{modality
        experts} (BERT, ResNet, Wav2Vec, etc.) feed \textbf{projection heads} into a
        shared space. This is a valid and common industrial pattern (compose experts,
        do not force one backbone).
  \item \textbf{Fusion vs retrieval:} the RAG path fuses information at the
        \textbf{token/context level} inside the LLM after retrieval, not only at the
        embedding level. Pipeline~E's Transformer fusion is a \emph{different}
        experimental head for MOSI-style sentiment; it should not be described as
        identical to the retrieval encoder stack without integration proof.
  \item \textbf{Domain transfer claim:} the radiology transfer story (retrain mainly
        the image adapter, freeze most parameters) is coherent with parameter-efficient
        adaptation practice. Any document that simultaneously claims ``only the image
        head'' and ``TextProjection minor nudge'' should pick one \emph{per experiment
        configuration} and report both if both were run.
  \item \textbf{Pipeline~C lesson:} joint InfoNCE between unrelated modalities without
        true paired semantics can \textbf{damage} the shared space; the standalone
        audio track is the correct engineering response. Proposals should cite this as
        risk mitigation, not hide it.
  \item \textbf{Pipeline~E limitation:} pseudo-labels add noise; accuracy gaps vs
        published supervised methods should be expected. The proposal must not imply
        parity with fully supervised SOTA without that caveat.
\end{itemize}

\section{Revised high-level architecture narrative}

Figure~\ref{fig:concept} summarises the \textbf{conceptual} architecture (components
may reside in separate services or notebooks in the prototype).

\begin{figure}[htbp]
\centering
\fbox{\parbox{0.92\textwidth}{
\small
\textbf{User} $\rightarrow$ (speech $\rightarrow$ ASR/embeddings; optional image)
$\rightarrow$ \textbf{Shared 512-d space}\\
$\rightarrow$ \textbf{FAISS} $\rightarrow$ \textbf{KG re-rank} $\rightarrow$
\textbf{LLM + citations/confidence} $\rightarrow$ (optional TTS)\\[0.5em]
\textbf{Parallel affect path:} audio (+\ optionally video frame) $\rightarrow$
frozen backbones $\rightarrow$ trained projections $\rightarrow$
\textbf{centroid or fusion head} $\rightarrow$ label/score for prompt modulation.
}}
\caption{Conceptual data flow: retrieval stack (top) and affect branch (bottom).
Training tracks A--E populate indices, projections, centroids, or fusion weights.}
\label{fig:concept}
\end{figure}

\section{Work packages (revised)}

\begin{enumerate}
  \item \textbf{Shared embedding API:} contract for unit-norm vectors, batching, and
        device placement; single source of truth for dimension $512$.
  \item \textbf{Index lifecycle:} build/rebuild FAISS for each domain; version metadata
        with corpus checksums.
  \item \textbf{KG + NER:} spaCy/scispaCy pipelines per domain; PMI edges; bounded DFS
        for neighbour expansion.
  \item \textbf{Generator integration:} prompt templates, context packing, refusal
        tiers, optional Groq/Llama backends as documented.
  \item \textbf{Affect modules:} export centroid tensors and projection checkpoints;
        document warm-start chain C$\to$D and separate evaluation splits (video-level
        for MOSI).
  \item \textbf{Evaluation:} retrieval metrics (P@$k$, MRR) on held-out captions;
        emotion per-class accuracy; for MOSI, separate reporting for regression vs
        binary tasks under pseudo-label noise.
  \item \textbf{Reproducibility:} pinned dependencies, saved caches for heavy
        encoders, and explicit notebook$\leftrightarrow$document version mapping.
\end{enumerate}

\section{Risks and mitigations}

\begin{description}[style=nextline,font=\normalfont\bfseries]
  \item[Vocabulary collision.] Always say \textbf{runtime stage} vs
        \textbf{training track} when using ``phase'' or ``pipeline.''
  \item[Scope creep.] Pipeline~E is optional research; the capstone demo can be
        complete without it if time-constrained.
  \item[Label noise.] Flag pseudo-labels in E; do not merge MOSI metrics with CREMA-D
        emotion numbers without context.
  \item[Integration debt.] Retrieval + LLM + TTS + affect is multiple subsystems;
        define minimal integration slice (e.g., text QA + audio emotion only) first.
\end{description}

\section{Deliverables}

\begin{itemize}[nosep]
  \item LaTeX technical reports per track (as in repository).
  \item Trained projection weights, FAISS indices, KG serialisations, centroid files.
  \item Runnable notebooks or scripts with README describing which artifact each
        produces (note: notebooks referenced in existing \texttt{.tex} files must be
        present in the submission bundle for full reproducibility).
  \item Short demo narrative mapping user utterance $\rightarrow$ retrieval $\rightarrow$
        grounded answer $\rightarrow$ optional emotional tone of response.
\end{itemize}

\end{document}
